\section*{अध्याय \ref{ch:b1:active-margins} --- सक्रिय किनारे}

\begin{solution}{exo:b1:active-margins:1}
(क) $2\pi G\,\Delta\rho\,h = 2\pi\times\num{6.674e-11}\times1700\times3000 =
\qty{2.14e-3}{m.s^{-2}}$: गुरुत्व लगभग \qty{214}{\milli\gal} कम होता। (ख) \qty{3.5}{km} के लिए, $2\pi\times\num{6.674e-11}\times1700\times3500 = \qty{2.50e-3}{m.s^{-2}}$, यानी \qty{250}{\milli\gal}: यह प्रेक्षित न्यूनता के आकार का है, अतः खाई समस्थैतिक साम्य से बहुत दूर रोकी गई है। (ग) स्लैब का सूत्र मान लेता है कि लुप्त चट्टान हर ओर अनंत तक फैली है; केवल दसियों किलोमीटर चौड़ी खाई बिंदु के आसपास बहुत कम द्रव्यमान हटाती है, अतः अकेला उसका पानी \qty{250}{\milli\gal} से कम देता है। शेष न्यूनता गहराई में है: समूची प्लेट खाई के नीचे मुड़ती है, और उसकी \omterm{def:b1:layered-earth:layers}{भूपर्पटी}, जो \omterm{def:b1:layered-earth:layers}{प्रावार} से हल्की है, अपने आधार पर \omterm{def:b1:layered-earth:layers}{प्रावार} का स्थान लेती है।
\end{solution}

\begin{solution}{exo:b1:active-margins:2}
(क) $\qty{82}{mm.yr^{-1}} = \qty{82}{km.Myr^{-1}}$, अतः \qty{10}{Myr} में \qty{820}{km}। (ख) $820\sin\ang{45} = \qty{580}{km}$। (ग) \qty{600}{km} तक का पथ $600/\sin\ang{45} = \qty{849}{km}$ है, जो $849/82 = \qty{10.3}{Myr}$ में तय होता है।
\end{solution}

\begin{solution}{exo:b1:active-margins:3}
$g = 110/\tan\delta$: \ang{30} के लिए \qty{191}{km}, \ang{45} के लिए \qty{110}{km}, \ang{60} के लिए \qty{64}{km}। मंद \omterm{def:b1:geological-maps:attitude}{नति} वाले स्लैब के ऊपर का चाप अपनी खाई से सबसे दूर पड़ता है।
\end{solution}

\begin{solution}{exo:b1:active-margins:4}
(क) अभिसरण के लिए: $\qty{2.7e6}{m^2.yr^{-1}}/\qty{5.1e7}{m} =
\qty{0.053}{m.yr^{-1}}$, लगभग \qty{53}{mm.yr^{-1}}। प्रसार के लिए: $\qty{3.0e6}{m^2.yr^{-1}}/\qty{6.7e7}{m} = \qty{0.045}{m.yr^{-1}}$, लगभग \qty{45}{mm.yr^{-1}}। एक छोटी कुल लंबाई पर सबडक्शन सीमाएँ महासागर-तल को कटक उसे बनाने से तेज़ नष्ट करती हैं, फिर भी कुल मिलाकर कम नष्ट करती हैं। (ख) शेष लगभग \qty{0.3}{km^2} मॉडल की अन्य अभिसारी सीमाओं पर खप जाता है, जैसे महाद्वीपीय टकराव और महासागरों के भीतर के अभिसारी क्षेत्र, जहाँ प्लेटें नष्ट होने के बजाय सिकुड़ती हैं; और कटक ही अपसारी सीमाएँ नहीं हैं: महाद्वीपीय दरारें और कुछ रूपांतर भी थोड़े फैलते हैं।
\end{solution}

\begin{solution}{exo:b1:active-margins:5}
(क) $\bar x = 350$, $\bar z = 290$; $x$ के विचलन $-150$, $-50$, $50$, $150$ हैं और $z$ के $-140$, $-40$, $40$, $140$, अतः $S_{xx} = 2(150^2 + 50^2) = \num{50000}$ और $S_{xz} = 2(150\times140 + 50\times40) = \num{46000}$। तब $s = 0.92$, $a = 290 - 0.92\times350 = \qty{-32}{km}$, और \omterm{def:b1:geological-maps:attitude}{नति} $\arctan 0.92 = \ang{42.6}$ है। (ख) $z = 0$ पर $x = 32/0.92 = \qty{34.8}{km}$, खाई से दूर। (ग) अंत-बिंदुओं से होकर ढलान $280/300 = 0.933$ है, \omterm{def:b1:geological-maps:attitude}{नति} \ang{43.0}: निकट, पर वह चार भूकंपों में से केवल दो का प्रयोग करती है।
\end{solution}

\begin{solution}{exo:b1:active-margins:6}
(क) $50/\tan\ang{25} + 60/\tan\ang{45} = 107.2 + 60.0 = \qty{167.2}{km}$। (ख) पूरे रास्ते \ang{45} का स्लैब \qty{110}{km} देता: मंद उथला खंड अंतर को \qty{57}{km} चौड़ा कर देता है। (ग) गहरा खंड $z = 50 + (x - 107.2)$ है; वह $z = 0$ पर $x = \qty{57.2}{km}$ तक पहुँचता है, खाई से भू-ओर, जैसा कुरील रेखा \qty{53}{km} पर करती है: एक मंद उथला खंड कुरील विस्थापन की व्याख्या करता है। (घ) $107.2 + 60/\tan\ang{60} = 107.2 + 34.6 =
\qty{141.9}{km}$।
\end{solution}

\begin{solution}{exo:b1:active-margins:7}
(क) त्रिज्या \omterm{def:b1:geological-maps:attitude}{नति} का आधा है, और $r = R\,\theta$ जहाँ $\theta$ रेडियन में है: \ang{45} के लिए \ang{22.5} और $6371\times0.3927 = \qty{2502}{km}$; \ang{30} के लिए \ang{15} और \qty{1668}{km}; \ang{70} के लिए \ang{35} और \qty{3892}{km}। (ख) मंद \omterm{def:b1:geological-maps:attitude}{नति} वाले स्लैबों के ऊपर के चाप। (ग) पूर्वानुमान $44.9/2 = \ang{22.45}$ है; फ़िट, \ang{21.9}, $0.55/22.45 = \qty{2.4}{\percent}$ छोटा है।
\end{solution}

\begin{solution}{exo:b1:active-margins:8}
(क) $q/k$: $0.035/2.5 = \qty{0.014}{K.m^{-1}} = \qty{14}{K.km^{-1}}$ अग्र भाग के नीचे, $0.080/2.5 = \qty{32}{K.km^{-1}}$ बैक-आर्क के नीचे। (ख) $10 + 14\times30 = \qty{430}{\celsius}$ और $10 + 32\times30 = \qty{970}{\celsius}$। (ग) सतह के \omterm{def:b1:heat-budget:heat-flow}{ऊष्मा-प्रवाह} का एक भाग \omterm{def:b1:layered-earth:layers}{भूपर्पटी} के भीतर की रेडियोधर्मिता से आता है, अतः गहरे स्तरों से कम ऊष्मा गुज़रती है और प्रवणता गहराई के साथ घटती है: सतह की प्रवणता का रैखिक बहिर्वेशन गहराई पर तापमान को अधिक आँकता है।
\end{solution}

\begin{solution}{exo:b1:active-margins:9}
(क) \qty{32}{mm.yr^{-1}} = \qty{32}{km.Myr^{-1}}; अर्ध दर \qty{16}{mm.yr^{-1}}। (ख) $100/32 = \qty{3.1}{Myr}$। (ग) दर का स्थिर होना आवश्यक नहीं; चौड़ाई का कुछ भाग प्रसार आरंभ होने से पहले चाप के खिंचने और दरार पड़ने से आ सकता है; और अक्ष कूद गया हो सकता है या असममित रूप से फैला हो।
\end{solution}

\begin{solution}{exo:b1:active-margins:10}
(क) पथ $600/\sin\ang{45} = \qty{849}{km}$ है, जो \qty{82}{km.Myr^{-1}} पर \qty{10.3}{Myr} में तय होता है, अर्थात् \qty{3.27e14}{s}। (ख) $\sqrt{\num{e-6}\times\num{3.27e14}} = \qty{1.8e4}{m}$, लगभग \qty{18}{km}। (ग) $(\qty{5e4}{m})^2/\qty{e-6}{m^2.s^{-1}} = \qty{2.5e15}{s}$, लगभग \qty{79}{Myr}। (घ) ऊष्मा स्लैब के हर पार्श्व के केवल बाहरी \qty{18}{km} तक पहुँची है: उसका केंद्र \qty{600}{km} पर लगभग उतना ही ठंडा पहुँचता है जितना खाई से चला था, अपने आसपास के \omterm{def:b1:layered-earth:layers}{प्रावार} से ठंडा और मज़बूत, और इसी कारण वह उस गहराई पर अब भी भूकंपों में टूट सकता है।
\end{solution}

\begin{solution}{exo:b1:active-margins:11}
(क) पैर के पीछे दूरी $L$ पर सतह पैर के स्तर से $L\tan\alpha$ ऊपर खड़ी है और आधार उसके नीचे $L\tan\beta$ पर है: पश्च-भाग $L(\tan\alpha + \tan\beta)$ ऊँचा है, और त्रिभुज का क्षेत्रफल $\frac12 L\times L(\tan\alpha + \tan\beta)$ है। (ख) $\tan\ang{4} + \tan\ang{6} =
0.175$। प्रिज़्म को $\qty{1}{km}\times\qty{40}{km.Myr^{-1}} =
\qty{40}{km^2.Myr^{-1}}$ खंड मिलता है, अतः $A = 40t$ और $L = \sqrt{2A/0.175}$: \qty{1}{Myr} के बाद \qty{21.4}{km}, \qty{4}{Myr} के बाद \qty{42.8}{km}, \qty{10}{Myr} के बाद \qty{67.6}{km}। $L$ $\sqrt t$ के अनुसार बढ़ता है: चार गुना समय दोगुनी चौड़ाई देता है। (ग) $A$ आधा होने पर $L$ $\sqrt2$ से भाग देने पर घटता है: \qty{15.1}{km}, \qty{30.2}{km} और \qty{47.8}{km}।
\end{solution}

\begin{solution}{exo:b1:active-margins:12}
(क) $\partial F/\partial a = -2\sum(z_i - a - s x_i) = 0$ देता है $n\bar z = na + sn\bar x$, अर्थात् $\bar z = a + s\bar x$। (ख) हर $z_i$ और $\bar z$ \qty{10}{km} बढ़ते हैं, अतः विचलन $z_i - \bar z$ और $S_{xz}$ अपरिवर्तित हैं: ढलान वही है, और $a = \bar z - s\bar x$ \qty{10}{km} बढ़ता है। (ग) विचलन $x_i - \bar x$ अपरिवर्तित हैं, अतः ढलान वही है; $\bar x$ 20 घटता है, अतः $a$ $20s$ बढ़ता है। (घ) कोई नहीं: \omterm{def:b1:geological-maps:attitude}{नति} केवल ढलान पर निर्भर करती है। सभी बिंदुओं में साझा त्रुटियाँ रेखा को हिलाती हैं, उसे घुमाती नहीं।
\end{solution}

\begin{solution}{pb:b1:active-margins:1}
\textbf{1.} आठों बिंदु लगभग \ang{45} \omterm{def:b1:geological-maps:attitude}{नति} वाली एक सीधी रेखा के निकट पड़ते हैं, खाई से \qty{155}{km} पर \qty{96}{km} गहरे से लेकर \qty{609}{km} पर \qty{573}{km} गहरे तक।

\textbf{2.} $\sum x_i = 3152$, $\sum z_i = 2751$, $\sum x_i^2 = \num{1442910}$, $\sum x_i z_i = \num{1290548}$; $\bar x = \qty{394.0}{km}$, $\bar z = \qty{343.875}{km}$।

\textbf{3.} $S_{xx} = 1442910 - 8\times394^2 = \num{201022}$ और $S_{xz} = 1290548 - 8\times394\times343.875 = \num{206654}$। तब $s = 206654/201022 = 1.028$ और $a = 343.875 - 1.02802\times394 =
\qty{-61.2}{km}$: $z = -61.2 + 1.028\,x$।

\textbf{4.} $\arctan1.028 = \ang{45.8}$।

\textbf{5.} खाई से भू-ओर $x = 61.2/1.028 = \qty{59.5}{km}$ पर। उस \omterm{def:b1:geological-maps:attitude}{नति} पर खाई से निकला स्लैब वहाँ सतह तक पहुँचता; वह सतह तक खाई पर ही पहुँचता है, इसलिए वह खाई के निकट अधिक मंद \omterm{def:b1:geological-maps:attitude}{नति} से झुकता है और नीचे तीखा होता है।

\textbf{6.} विचलन $-2.1$, $2.0$, $-7.3$, $5.3$, $7.1$, $2.3$, $-15.0$ और $8.1$ \unit{km} हैं। भूकंप 7 रेखा से \qty{15}{km} ऊपर है, भूकंप 8 \qty{8}{km} नीचे; बाक़ी लगभग \qty{7}{km} के भीतर हैं, जो स्लैब की मोटाई का एक भाग है।

\textbf{7.} $\ang{320.9} - \ang{300} = \ang{20.9}$। खाई के लंबवत: $87.4\cos\ang{20.9} = \qty{81.6}{mm.yr^{-1}}$; समांतर: $87.4\sin\ang{20.9} = \qty{31.2}{mm.yr^{-1}}$। अभिसरण तिरछा है।

\textbf{8.} \qty{81.6}{mm.yr^{-1}} = \qty{81.6}{km.Myr^{-1}}: \qty{10}{Myr} में \qty{816}{km}।

\textbf{9.} पथ $632/\sin\ang{45} = \qty{894}{km}$ है, जो $894/81.6 = \qty{11.0}{Myr}$ में तय होता है।

\textbf{10.} $122 + 11 = \qty{133}{Myr}$।

\textbf{11.} $t = \qty{10.95}{Myr} = \qty{3.46e14}{s}$ और $\sqrt{\num{e-6}\times\num{3.46e14}} = \qty{1.86e4}{m}$, हर पार्श्व से लगभग \qty{19}{km}: \qty{100}{km} मोटे स्लैब के मध्य की \qty{50}{km} दूरी से बहुत कम, जिसका केंद्र अब भी ठंडा है।

\textbf{12.} $q/k$: पुराने महासागर-तल के नीचे \qty{17.2}{K.km^{-1}}, अग्र भाग के नीचे \qty{14.0}{K.km^{-1}} और बैक-आर्क के नीचे \qty{32.0}{K.km^{-1}}।

\textbf{13.} अग्र भाग के नीचे ठंडा स्लैब केवल कुछ दसियों किलोमीटर नीचे है; वह अपने ऊपर की चट्टान से ऊष्मा लेता है, अतः सतह तक कम ऊष्मा पहुँचती है।

\textbf{14.} $h = \qty{2.60e-3}{m.s^{-2}}/(2\pi\times\num{6.674e-11}\times1700)
= \qty{3647}{m}$, लगभग \qty{3.6}{km}: उसके आसपास के तल से नीचे खाई की पूरी गहराई, लगभग \qty{3.5}{km}, जितना, या उससे कुछ अधिक। खाई से ऐसा द्रव्यमान गायब है जिसकी कोई क्षतिपूर्ति नहीं करता।

\textbf{15.} घनत्व अंतर $3300 - 1000 = \qty{2300}{kg.m^{-3}}$ है, और $2\pi\times\num{6.674e-11}\times2300\times3500 = \qty{3.38e-3}{m.s^{-2}}$, यानी \qty{338}{\milli\gal}। स्लैब का सूत्र लुप्त द्रव्यमान को किसी भी दूरी तक गिनता है; वास्तविक खाई केवल दसियों किलोमीटर चौड़ी पट्टी पर नीचे मुड़ी है, अतः उसकी न्यूनता छोटी है, लगभग \qty{260}{\milli\gal}, दोनों स्लैब-अनुमानों के बीच।

\textbf{16.} मुड़ती प्लेट की सामर्थ्य, जो खाई को नीचे और बाहरी उभार को ऊपर उसी तरह रोके रखती है जैसे एक प्रत्यास्थ प्लेट रखती, और सघन स्लैब का \omterm{def:b1:layered-earth:layers}{प्रावार} में धँसने का खिंचाव।

\textbf{17.} $-61.2 + 1.028\times164 = \qty{107.4}{km}$।

\textbf{18.} हाँ: \qty{107}{km}, 72 और \qty{173}{km} के बीच है, और औसत, \qty{105}{km}, से केवल लगभग \qty{2}{km} गहरा है।

\textbf{19.} $x = (100 + 61.2)/1.028 = \qty{156.8}{km}$।

\textbf{20.} $\ang{45.8}/2 = \ang{22.9}$, और $6371\times22.9\times\pi/180 = \qty{2546}{km}$।

\textbf{21.} फ़िट, \ang{21.9}, \ang{1.0} छोटा है, लगभग \qty{4}{\percent}: मुड़ा कवच कुरील चाप की वक्रता को अच्छी तरह वर्णित करता है।

\textbf{22.} $x = (105 + 61.2)/1.028 = 166.2/1.028 = \qty{161.7}{km}$।

\textbf{23.} $164 - 161.7 = \qty{2.3}{km}$, प्रेक्षित चाप-मोर्चे से खाई के अधिक निकट, $2.3/164 = \qty{1.4}{\percent}$: कुरील ज्वालामुखी ऐसे स्लैब के ऊपर खड़े हैं जो औसत से बहुत थोड़ा गहरा है, जैसा प्रश्न 18 में पाया गया।

\textbf{24.} निकटतम किलोमीटर तक, कुरील स्लैब रेखा \qty{105}{km} पर स्लैब के ऊपर के ज्वालामुखियों के लिए \omterm{def:b1:active-margins:volcanic-arc}{चाप–खाई अंतराल} का अनुमान \textbf{\qty[text-series-to-math=true]{162}{km}} लगाती है (प्रेक्षित: \qty{164}{km})।
\end{solution}
